%% ====================================================================
%%  Chapter 1 — Foundations: WiFi Localization and the CSI Signal
%% ====================================================================
\chapter{Foundations: WiFi Localization and the CSI Signal}
\label{ch:foundations}

Before any neural network appears, you need to understand the physical problem and the clever
reframing that makes deep learning applicable to it. This chapter builds that foundation: why indoor
localization is hard, what a WiFi signal actually measures, how that measurement becomes an
``image,'' and why the whole problem is then solved as \emph{image-to-image translation}.

\section{Why indoor localization, and why GPS fails indoors}
\label{sec:found:gps}

Many applications need to know \emph{where} something is inside a building: tracking assets in a
warehouse, finding patients in a hospital, augmented reality, indoor robots, and location-aware
services in malls or airports. Outdoors, GPS solves this to within a few meters. Indoors, it
collapses.

GPS works by timing signals from satellites in direct line of sight. Indoors, those signals are
\emph{attenuated} (weakened) by walls and roofs and \emph{reflected} off surfaces, so a receiver sees
a jumble of delayed copies rather than a clean direct path. The timing that GPS depends on is
destroyed. Two decades of research therefore looked for indoor alternatives: WiFi, Bluetooth,
ultra-wideband, and inertial sensors.

\begin{keyidea}
The same physics that breaks GPS indoors ---reflection and \emph{multipath}--- is exactly what an
indoor WiFi system must learn to cope with. Our models do not avoid multipath; they learn to read
through it.
\end{keyidea}

\section{The WiFi opportunity: access points and CSI}
\label{sec:found:wifi}

WiFi is the most studied indoor-localization technology for one practical reason: the infrastructure
already exists. Almost every building is blanketed by WiFi access points (APs). And modern WiFi
hardware, as a side effect of doing its real job (reliable communication), continuously measures
something extremely useful for localization: \textbf{Channel State Information (CSI)}.

When an AP communicates with a device, it must estimate how the wireless channel distorted the
signal so it can undo that distortion. That per-channel estimate \emph{is} the CSI. Because the way a
signal is distorted depends on the geometry between the device and the AP ---distance, angle, what it
bounced off--- the CSI implicitly encodes \emph{where the device is}. Extracting that location is the
whole game.

\section{What CSI physically is}
\label{sec:found:csi}

WiFi uses \emph{orthogonal frequency-division multiplexing} (OFDM): instead of one carrier frequency,
the signal is split across many narrow \textbf{subcarriers}. For each subcarrier (and each antenna on
the AP), the hardware estimates a single complex number describing how that frequency was attenuated
and phase-shifted by the channel. Stack those numbers up and you get the CSI matrix at AP~$i$:
\[
  H_i \in \mathbb{C}^{\,N_{\text{sub}} \times N_{\text{ant}}},
\]
with $N_{\text{sub}}$ subcarriers and $N_{\text{ant}}$ antennas. Each entry has an \emph{amplitude}
(how much the signal weakened) and a \emph{phase} (how much it was delayed, modulo a wavelength).

Two facts make CSI informative about location:
\begin{itemize}[leftmargin=1.4em,nosep]
  \item \textbf{Amplitude vs.\ distance.} Signal strength falls off with distance, so amplitude is a
  (noisy) distance cue.
  \item \textbf{Phase across antennas vs.\ angle.} A wavefront arriving at an angle hits the AP's
  antennas at slightly different times, producing a predictable phase progression across the antenna
  array. That progression encodes the \emph{angle of arrival}.
\end{itemize}

\section{From CSI to a per-AP heatmap: AoA and ToF}
\label{sec:found:aoa}

Raw complex CSI is not something you would feed a vision network directly. DLoc's pre-processing
(done in MATLAB, ``stage 0'' of the pipeline) turns each AP's CSI into a 2D \textbf{Angle-of-Arrival
/ Time-of-Flight (AoA--ToF) profile} and then projects it onto the physical floor grid. The result is
a per-AP \emph{heatmap} over the room: bright where ``the signal looks like it came from here.''

\begin{itemize}[leftmargin=1.4em,nosep]
  \item \textbf{Angle of Arrival (AoA)} ---from the phase progression across antennas--- says
  \emph{which direction} the signal came from.
  \item \textbf{Time of Flight (ToF)} ---from the phase progression across subcarriers (frequencies)---
  says \emph{how far} the signal travelled.
\end{itemize}

Combine direction and distance for one AP and you get a constraint on where the device is; combine
the constraints from several APs and you can, in principle, pin down the location. For the Jacobs
Hall environment used throughout the thesis there are $N_{\text{AP}}=4$ APs, and each heatmap is
sampled on a $161 \times 361$ grid at $0.1$\,m per pixel ---i.e.\ a $16.1 \times 36.1$\,m floor.
Stacking the four per-AP heatmaps gives a four-channel ``image'' of shape $[\,4,\,161,\,361\,]$.

\begin{figure}[H]
\centering
\begin{tikzpicture}[
  font=\small,
  stage/.style={draw=accent,fill=accent!8,rounded corners,minimum height=1.05cm,
    minimum width=2.5cm,align=center},
  arr/.style={-{Stealth[length=2.4mm]},accent,thick}]
\node[stage] (csi) {Raw CSI\\$H_i\in\mathbb{C}^{N_\text{sub}\times N_\text{ant}}$\\(4 APs)};
\node[stage,right=0.9cm of csi] (prof) {AoA--ToF\\profiles\\(MATLAB)};
\node[stage,right=0.9cm of prof] (img) {4-channel\\input image\\$[4,161,361]$};
\node[stage,right=0.9cm of img,fill=accent2!12,draw=accent2!70!black] (net) {Neural\\network\\(translate)};
\node[stage,right=0.9cm of net] (out) {Output\\heatmap\\$[1,161,361]$};
\node[stage,below=0.7cm of out] (xy) {$\arg\max$\\$\Rightarrow (x,y)$};
\draw[arr] (csi) -- (prof);
\draw[arr] (prof) -- (img);
\draw[arr] (img) -- (net);
\draw[arr] (net) -- (out);
\draw[arr] (out) -- (xy);
\end{tikzpicture}
\caption{The end-to-end pipeline. Physics and signal processing turn raw CSI into a four-channel
``image'' (one channel per AP); the neural network translates that into a single-channel location
heatmap; the brightest pixel, scaled by $0.1$\,m/pixel, is the predicted position.}
\label{fig:found:pipeline}
\end{figure}

\section{The clock-offset problem}
\label{sec:found:offset}

There is a crucial wrinkle. ToF is a \emph{timing} measurement, and the device's clock is not
synchronized with the AP's clock. An unknown clock offset adds an unknown delay to every ToF, which
shifts each AP's heatmap \emph{radially} ---the apparent location can move by $10$--$15$\,m. Worse,
each AP has its own offset, so the four per-AP heatmaps do \emph{not} agree on a single point. The
preprocessing therefore produces two versions of the per-AP heatmaps:
\begin{itemize}[leftmargin=1.4em,nosep]
  \item \code{features\_w\_offset}: \emph{with} the offset present ---this is the network input
  (it is all you would have at deployment);
  \item \code{features\_wo\_offset}: with the offset removed using extra information available only at
  training time ---this is a \emph{training target} that teaches the network to undo the offset.
\end{itemize}
This is the single most important reason a \emph{learned} system beats classical geometry here:
resolving the offset and reconciling four disagreeing APs is exactly what the network learns to do
(the DLoc ``consistency decoder''; see Ch.~\ref{ch:arch} and \code{docs/01}).

\section{Localization as image-to-image translation}
\label{sec:found:translation}

Now the reframing that makes everything click. Instead of computing the location with geometry, DLoc
treats the four-channel CSI image as the \emph{input} of an \textbf{image-to-image translation}
problem ---the same family of problems as ``turn this sketch into a photo''--- and learns a network
that outputs a clean, single-channel heatmap. The target heatmap is a 2D \emph{Gaussian bump}
centered on the true location (\code{labels\_gaussian\_2d}, shape $[1,161,361]$). At test time the
prediction is read off as the \textbf{argmax} of the output heatmap, scaled by $0.1$\,m/pixel; the
error is the Euclidean distance from that pixel to the true one (formalized in Ch.~\ref{ch:eval}).

\begin{keyidea}
Three reframings turn a hard physics problem into a tractable learning problem: (1) represent each
AP's signal as a \emph{spatial heatmap} over the floor; (2) make the network's job \emph{translate
input heatmaps into one output heatmap}; (3) read the location off as the \emph{peak}. Every model in
this thesis ---except Mamba--- accepts this image framing, which is precisely why Mamba struggles
(Ch.~\ref{ch:arch}).
\end{keyidea}

\section{Why deep learning instead of classical geometry}
\label{sec:found:why_dl}

Classical methods (e.g.\ SpotFi) estimate AoA and ToF analytically and intersect them. They work in
clean, line-of-sight conditions but degrade badly under \emph{non-line-of-sight} (NLOS) and rich
indoor multipath, where the strongest path is often a reflection, not the direct one. A learned
network can instead exploit statistical regularities across the four APs and across the training
environment: it can pick the direct path by cross-checking APs, sharpen blurry profiles, and undo
the clock offset. The price for this capability is model size ---and shrinking that price, without
losing the capability, is the subject of the rest of this guide.
